% !TEX program = xelatex
% SKED document (Socratic Knowledge Elicitation & Documentation), phase 5.
% Build: python3 <plugin>/scripts/build.py --dir .   (or: xelatex main; xelatex main)
% Every paragraph that carries knowledge from K* ends with \sked{K..} tags.
% Add the `review` class option to see those tags in the margin.
\documentclass[profile=dstu,lang=uk]{skedreport}

\skedtitle{Специфікація вимог до інформаційної системи персонального обліку медіа-контенту}
\skedsubtitle{Software Requirements Specification}
\skedauthor{Букша Вікторія}
\skedaffiliation{Студентка курсу «Проектування інформаційних систем»}
\skedorganization{}
\skeddoctype{Специфікація вимог}
\skedversion{1.3.0}
\skeddate{25 вересня 2026~р.}
\skedcity{Київ}
\skedudc{004.415.2}
\skedkeywords{специфікація вимог, медіа-трекер, веб-застосунок, SRS, каталог контенту, сократичний діалог}
\skedplatform{Веб-браузер}

\begin{document}

\skedmaketitle

\begin{abstract}
Цей документ містить специфікацію вимог (SRS) до веб-застосунку персонального обліку медіа-контенту~--- книг, серіалів і фільмів. Система призначена для зареєстрованих користувачів, які ведуть особистий каталог переглянутого та прочитаного контенту. Документ визначає функціональні вимоги (31~вимога), нефункціональні вимоги (8~вимог), критерії прийняття, BDD-сценарії та матрицю простежуваності для версії~v1 (MVP). Документ підготовлений за методологією SKED (Socratic Knowledge Elicitation \& Documentation) у рамках навчального курсу «Проектування інформаційних систем». До MVP входять: реєстрація та авторизація, каталог із CRUD-операціями, теги з автодоповненням, механізм Sessions для фіксації повторних переглядів, добірки (many-to-many), автозаповнення дати завершення та особиста статистика. Функції соціальної взаємодії, офлайн-режим та інтеграція зі сторонніми медіа-базами перенесені до наступних версій.
\end{abstract}

\skedtoc

\section{Вступ}

\subsection{Мета та область застосування документа}

Цей документ є специфікацією програмних вимог (Software Requirements Specification, SRS) до інформаційної системи персонального обліку медіа-контенту. Мета документа~--- формально описати функціональні та нефункціональні вимоги до системи, визначити межі версії~v1 (MVP) та надати базу для розроблення проектної документації (SDD) і тестування. \sked{K10}

\subsection{Призначення системи}

Система є веб-застосунком, що дозволяє зареєстрованому користувачеві вести особистий каталог медіа-контенту~--- книг, серіалів і фільмів. Застосунок забезпечує облік переглянутого та прочитаного контенту: фіксацію статусів, оцінок, нотаток, тегів, повторних переглядів (Sessions), організацію елементів у добірки (Collections) і перегляд особистої статистики. \sked{K2,K4,K3,K5}

\subsection{Контекст і аудиторія документа}

Документ підготовлений у рамках навчального курсу «Проектування інформаційних систем» (SDD-цикл). Цільова аудиторія: викладач курсу та розробники системи. Читач цього документа повинен отримати достатньо інформації для розроблення архітектури та плану тестування без додаткових уточнень. \sked{K44}

\subsection{Структура документа}

\begin{itemize}
  \item \secref{sec:zahalnyi-opys}~--- загальний опис системи, межі MVP, ролі користувачів, операційне середовище.
  \item \secref{sec:hlosarii}~--- визначення ключових термінів.
  \item \secref{sec:funktsionalni-vymohy}~--- повний перелік функціональних вимог, згрупованих за функціональними блоками.
  \item \secref{sec:nefunktsionalni-vymohy}~--- нефункціональні вимоги: продуктивність, безпека, надійність, зручність, сумісність, валідація.
  \item \secref{sec:kryterii-pryiniattia}~--- критерії прийняття для кожної вимоги.
  \item \secref{sec:bdd}~--- BDD-сценарії у форматі Gherkin.
  \item \secref{sec:matrytsia}~--- матриця простежуваності вимог.
  \item Додаток~А~--- перелік функцій, що не входять до MVP.
\end{itemize}

\section{Загальний опис системи}
\label{sec:zahalnyi-opys}

\subsection{Загальна характеристика та перспектива продукту}

Система є самостійним веб-застосунком без інтеграцій із зовнішніми сервісами в межах MVP. Кожен зареєстрований користувач має власний ізольований простір даних~--- каталог, добірки та статистику. Система не є частиною більшої платформи і не надає публічного API у версії~v1. \sked{K38,K41,K42}

\subsection{Склад функціональних блоків}

Система реалізує такі функціональні блоки:
\begin{enumerate}
  \item Автентифікація та управління акаунтом.
  \item Особистий каталог медіа-контенту.
  \item Управління елементами контенту (CRUD).
  \item Теги та фільтрування за тегами.
  \item Sessions~--- фіксація повторних переглядів або прочитань.
  \item Добірки~--- власні іменовані групи елементів.
  \item Автозаповнення дати завершення.
  \item Особиста статистика.
\end{enumerate}
\sked{K10,K13,K16,K18,K21,K27,K28}

\subsection{Ролі користувачів}

У версії~v1 існує єдина роль. Адміністративна роль відсутня.

\begin{table}[h]
\caption{Ролі користувачів системи}
\label{tab:roli}
\begin{tabular}{L{3cm} L{10cm}}
\toprule
\textbf{Роль} & \textbf{Опис} \\
\midrule
Користувач (User) & Зареєстрований власник акаунту. Має доступ виключно до власного каталогу, добірок і статистики. Може виконувати всі CRUD-операції над власними даними. \\
\bottomrule
\end{tabular}
\end{table}
\sked{K44}

\subsection{Операційне середовище}

Система реалізується як веб-застосунок. Нативний мобільний застосунок у версії~v1 не передбачений. Система повинна коректно функціонувати в останніх двох версіях браузерів Chrome, Firefox, Edge та Safari, включно з мобільним Safari. Мова реалізації бекенду не регламентується цим документом і визначається на етапі проектування архітектури. \sked{K38,K39,K45}

\subsection{Обмеження проектування та реалізації}

\begin{itemize}
  \item Реалізація виключно у вигляді веб-застосунку.
  \item У версії~v1 підтримується лише один мовний інтерфейс.
  \item Система не інтегрується із зовнішніми базами медіа-контенту (TMDB, Google Books тощо).
  \item Кожен користувач має доступ виключно до власних даних; крос-користувацький доступ не передбачений.
\end{itemize}
\sked{K38,K40,K41,K42}

\subsection{Межі MVP~--- що не входить до версії~v1}

Наступні функції є технічно можливими, але свідомо виключені з MVP для стримання обсягу першої версії. Детальний перелік наведено у Додатку~А. \sked{K43,K50}

\section{Глосарій}
\label{sec:hlosarii}

У \tabref{tab:hlosarii} наведено визначення ключових термінів, використаних у цьому документі.

\begin{table}[h]
\caption{Глосарій системи}
\label{tab:hlosarii}
\begin{tabular}{L{3.5cm} L{3.5cm} L{7cm}}
\toprule
\textbf{Термін (укр.)} & \textbf{Англ. відповідник} & \textbf{Визначення} \\
\midrule
Каталог & Catalog & Повний список усіх елементів контенту, доданих користувачем. \sked{K1} \\
\addlinespace
Елемент контенту & Content Item & Одна позиція в каталозі~--- книга, серіал або фільм із власними атрибутами (тип, статус, оцінка, теги тощо). \sked{K2} \\
\addlinespace
Добірка & Collection & Створена користувачем іменована група довільних елементів каталогу; один елемент може входити до кількох добірок одночасно. \sked{K3} \\
\addlinespace
Session & Session & Окремий запис факту перегляду або прочитання елемента; містить дату, оцінку та нотатку; для одного елемента може бути кілька Sessions. \sked{K4} \\
\addlinespace
Тег & Tag & Довільна текстова мітка, яку користувач додає до елемента для категоризації. \sked{K5} \\
\addlinespace
Статус & Status & Поточний стан роботи з елементом: \emph{заплановано~/ у процесі~/ завершено~/ відкладено}. \sked{K6} \\
\bottomrule
\end{tabular}
\end{table}

\section{Методологія розробки}
\label{sec:metodolohiia}

\subsection{Розробка на основі специфікації (SDD)}

Система розробляється відповідно до методології Specification Driven Development (SDD). Специфікація є контролюючим джерелом поведінки системи~--- жодний вихідний код не може визначати поведінку незалежно від цього документа. \sked{K53}

Обов'язковий процес при кожній зміні поведінки системи:

\begin{enumerate}
  \item Оновити цей документ до внесення змін до коду.
  \item Призначити або оновити ідентифікатор вимоги (REQ-F-NNN або REQ-NF-NNN).
  \item Визначити критерій прийняття (AC-NNN) для кожної вимоги.
  \item Написати автоматизований тест до реалізації (TDD, \secref{sec:metodolohiia}).
  \item Запустити тест і переконатися, що він провалюється з очікуваної причини (Red).
  \item Реалізувати мінімальний код, що проходить тест (Green).
  \item Рефакторити лише після успішного проходження тестів (Refactor).
  \item Оновити log-артефакти після завершення ітерації.
\end{enumerate}

\subsection{Обов'язкова розробка через тестування (TDD)}

TDD є обов'язковою для усієї реалізаційної роботи. Цикл Red~/ Green~/ Refactor застосовується до кожної функціональної зміни: \sked{K53}

\begin{itemize}
  \item \textbf{Red}~--- спочатку написати тест, що провалюється.
  \item \textbf{Green}~--- реалізувати мінімальний код, що проходить тест.
  \item \textbf{Refactor}~--- покращити код без зміни зовнішньої поведінки.
\end{itemize}

Жоден виробничий код не приймається, якщо він не покритий принаймні одним автоматизованим тестом. Реальні зовнішні системи (СУБД, сторонні API) не використовуються в unit-тестах~--- застосовуються детерміновані замінники (stubs / mocks).

\subsection{Правило трасованості}

Кожна реалізована поведінка системи повинна мати повну трасованість у ланцюжку:

\begin{codelisting}
\caption{Ланцюжок трасованості}
\label{lst:trace}
\begin{skedcode}
REQ → AC → Test → Implementation
\end{skedcode}
\end{codelisting}

Функціональність без ідентифікатора вимоги не є частиною системи. Вимога без критерію прийняття не готова до реалізації.

\subsection{Log-артефакти тестування}

Після кожного запуску тестів агент-розробник записує log-файл із результатами. \sked{K54} Log-файл повинен містити:

\begin{itemize}
  \item загальну кількість тестів і кількість passed~/ failed;
  \item статус кожного окремого тесту (passed~/ failed~/ skipped);
  \item час виконання тестового набору.
\end{itemize}

Log-файл є машиночитаним артефактом верифікації~--- на його основі агент-розробник і людина-рецензент визначають, чи виконане завдання відповідно до специфікації (крок «Is~verified?» у SDLC-процесі).

\section{Функціональні вимоги}
\label{sec:funktsionalni-vymohy}

\subsection{Автентифікація та управління акаунтом}

\begin{table}[h]
\caption{Вимоги до автентифікації та акаунту}
\label{tab:auth}
\begin{tabular}{L{2.5cm} L{11.5cm}}
\toprule
\textbf{ID} & \textbf{Вимога} \\
\midrule
REQ-F-001 & Система повинна надавати можливість реєстрації за email, паролем та нікнеймом (усі три поля~--- обов'язкові). Повторна реєстрація з тим самим email неможлива. \sked{K7} \\
\addlinespace
REQ-F-002 & Система повинна надавати вхід до акаунту за email та паролем. \sked{K8} \\
\addlinespace
REQ-F-003 & Система повинна надавати можливість виходу з акаунту. Після виходу доступ до захищених сторінок вимагає повторної авторизації. \sked{K9} \\
\addlinespace
REQ-F-029 & Користувач може змінити нікнейм у налаштуваннях акаунту. \sked{K49} \\
\bottomrule
\end{tabular}
\end{table}

\subsection{Каталог}

\begin{table}[h]
\caption{Вимоги до каталогу}
\label{tab:kataloh}
\begin{tabular}{L{2.5cm} L{11.5cm}}
\toprule
\textbf{ID} & \textbf{Вимога} \\
\midrule
REQ-F-004 & Система повинна відображати особистий каталог~--- список усіх елементів контенту користувача. \sked{K10} \\
\addlinespace
REQ-F-005 & Система повинна надавати фільтрування каталогу за типом елемента (книга~/ серіал~/ фільм). \sked{K11} \\
\addlinespace
REQ-F-006 & Система повинна надавати фільтрування каталогу за статусом елемента. Доступні значення статусу: «заплановано», «у процесі», «завершено», «відкладено». \sked{K12} \\
\addlinespace
REQ-F-023 & Система повинна надавати пошук елементів каталогу за назвою. \sked{K29} \\
\addlinespace
REQ-F-026 & Система повинна надавати сортування каталогу: за датою додавання (новіші~/ старіші) та за оцінкою (від вищої до нижчої і навпаки). \sked{K32} \\
\bottomrule
\end{tabular}
\end{table}

\subsection{Управління елементами контенту}

\begin{table}[h]
\caption{Вимоги до управління елементами контенту}
\label{tab:content}
\begin{tabular}{L{2.5cm} L{11.5cm}}
\toprule
\textbf{ID} & \textbf{Вимога} \\
\midrule
REQ-F-007 & Система повинна надавати можливість додавання нового елемента контенту з атрибутами: назва (обов'язково), тип (обов'язково), статус, оцінка~(0–10), нотатки, теги. Форма відображає тип-специфічні необов'язкові поля (REQ-F-024) та поле обкладинки (REQ-F-025). \sked{K13} \\
\addlinespace
REQ-F-008 & Система повинна надавати можливість редагування будь-якого атрибута наявного елемента контенту. \sked{K14} \\
\addlinespace
REQ-F-009 & Система повинна надавати можливість видалення елемента контенту з каталогу. Видалений елемент автоматично зникає з усіх добірок. \sked{K15} \\
\addlinespace
REQ-F-024 & Форма додавання та редагування елемента повинна відображати тип-специфічні необов'язкові поля залежно від обраного типу (см.~\tabref{tab:type-fields}). \sked{K30} \\
\addlinespace
REQ-F-025 & Кожен елемент контенту повинен мати обкладинку: за замовчуванням~--- одна з трьох ілюстрацій відповідно до типу. Користувач може замінити її власним зображенням (необов'язково). \sked{K31} \\
\addlinespace
REQ-F-028 & Формати зображення для обкладинки~--- JPEG, PNG, WebP; максимальний розмір файлу~--- 5~МБ. Система повідомляє користувача при порушенні обмежень. \sked{K48} \\
\bottomrule
\end{tabular}
\end{table}

\begin{table}[h]
\caption{Тип-специфічні поля елементів контенту}
\label{tab:type-fields}
\begin{tabular}{L{3cm} L{11cm}}
\toprule
\textbf{Тип} & \textbf{Необов'язкові поля} \\
\midrule
Книга & Автор, кількість сторінок \\
Серіал & Кількість сезонів \\
Фільм & Режисер, тривалість \\
\bottomrule
\end{tabular}
\end{table}
\sked{K30}

\subsection{Теги}

\begin{table}[h]
\caption{Вимоги до тегів}
\label{tab:tehy}
\begin{tabular}{L{2.5cm} L{11.5cm}}
\toprule
\textbf{ID} & \textbf{Вимога} \\
\midrule
REQ-F-010 & Система повинна надавати можливість додавання довільних тегів до елемента контенту. При введенні тегу система пропонує автодоповнення на основі раніше введених тегів користувача. Один елемент може мати кілька тегів. Тег можна видалити в будь-який момент. \sked{K16} \\
\addlinespace
REQ-F-011 & Система повинна надавати фільтрування та пошук по каталогу за тегом. \sked{K17} \\
\bottomrule
\end{tabular}
\end{table}

\subsection{Sessions (повторні перегляди~/ прочитання)}

\begin{table}[h]
\caption{Вимоги до Sessions}
\label{tab:sessions}
\begin{tabular}{L{2.5cm} L{11.5cm}}
\toprule
\textbf{ID} & \textbf{Вимога} \\
\midrule
REQ-F-012 & Система повинна надавати можливість фіксації кількох Sessions для одного елемента контенту; кожна Session має власні дату, оцінку та нотатку. \sked{K18} \\
\addlinespace
REQ-F-013 & Система повинна відображати список усіх Sessions елемента контенту у зворотному хронологічному порядку (від найновішої до найстарішої), із зазначенням загальної кількості. \sked{K19} \\
\addlinespace
REQ-F-014 & Система повинна надавати можливість редагування та видалення окремої Session. \sked{K20} \\
\bottomrule
\end{tabular}
\end{table}

\subsection{Добірки}

\begin{table}[h]
\caption{Вимоги до добірок}
\label{tab:dobirkы}
\begin{tabular}{L{2.5cm} L{11.5cm}}
\toprule
\textbf{ID} & \textbf{Вимога} \\
\midrule
REQ-F-015 & Система повинна надавати можливість додавання одного елемента каталогу до кількох добірок одночасно (зв'язок many-to-many). \sked{K21} \\
\addlinespace
REQ-F-016 & Система повинна надавати можливість створення нової іменованої добірки. \sked{K22} \\
\addlinespace
REQ-F-017 & Система повинна надавати можливість редагування назви добірки. \sked{K23} \\
\addlinespace
REQ-F-018 & Система повинна надавати можливість видалення добірки. Видалення добірки не видаляє елементи з каталогу. \sked{K24} \\
\addlinespace
REQ-F-019 & Система повинна надавати можливість видалення елемента з добірки без видалення його з каталогу чи з інших добірок. \sked{K25} \\
\addlinespace
REQ-F-020 & Система повинна відображати вміст добірки~--- список елементів, що до неї входять, з назвою, типом і статусом кожного. \sked{K26} \\
\addlinespace
REQ-F-027 & Система повинна надавати сортування елементів всередині добірки: за датою додавання (новіші~/ старіші) та за оцінкою (від вищої до нижчої і навпаки). \sked{K33} \\
\bottomrule
\end{tabular}
\end{table}

\subsection{Автозаповнення дати завершення}

\begin{table}[h]
\caption{Вимога до дати завершення}
\label{tab:data}
\begin{tabular}{L{2.5cm} L{11.5cm}}
\toprule
\textbf{ID} & \textbf{Вимога} \\
\midrule
REQ-F-021 & При зміні статусу елемента на «завершено» поле дати завершення повинно автоматично заповнюватися поточною датою, якщо дату не введено вручну. \sked{K27} \\
\bottomrule
\end{tabular}
\end{table}

\subsection{Особиста статистика}

\begin{table}[h]
\caption{Вимоги до особистої статистики}
\label{tab:stat}
\begin{tabular}{L{2.5cm} L{11.5cm}}
\toprule
\textbf{ID} & \textbf{Вимога} \\
\midrule
REQ-F-022 & Система повинна відображати особисту статистику: кількість елементів зі статусом «завершено» у розрізі типів (книги~/ серіали~/ фільми) та за обраний часовий проміжок. \sked{K28} \\
\addlinespace
REQ-F-022a & Вибір часового проміжку для статистики~--- з наперед визначених варіантів: «цей місяць», «цей рік», «весь час». \sked{K51} \\
\bottomrule
\end{tabular}
\end{table}

\section{Нефункціональні вимоги}
\label{sec:nefunktsionalni-vymohy}

\subsection{Продуктивність}

\textbf{REQ-NF-001.} Сторінки системи повинні завантажуватися в прийнятний для користувача час за нормальних умов мережевого підключення. \sked{K34}

\subsection{Безпека}

\textbf{REQ-NF-002.} Паролі користувачів зберігаються виключно у захищеному вигляді~--- у вигляді хешів (наприклад, bcrypt). Доступ до захищених ресурсів вимагає активної автентифікованої сесії. \sked{K35}

\textbf{REQ-NF-007.} Авторизаційний токен зберігається в HTTP-only cookie для запобігання XSS-атакам. Всі запити, що змінюють стан системи, повинні бути захищені від CSRF-атак. \sked{K47}

\subsection{Надійність}

\textbf{REQ-NF-003.} Усі дані користувача зберігаються в базі даних. Жодна дія користувача не повинна призводити до втрати раніше введених даних без явного підтвердження видалення з боку користувача. \sked{K36}

\subsection{Зручність використання}

\textbf{REQ-NF-004.} Інтерфейс системи повинен бути адаптивним для десктопних браузерів і мобільних браузерів. Окрема оптимізація під планшет у версії~v1 не передбачена. \sked{K37}

\textbf{REQ-NF-008.} При порожньому каталозі, порожній добірці або відсутніх результатах пошуку~/ фільтру система повинна відображати інформативне повідомлення замість порожнього списку. \sked{K52}

\subsection{Сумісність}

\textbf{REQ-NF-005.} Система повинна коректно функціонувати в останніх двох версіях браузерів: Google~Chrome, Mozilla~Firefox, Microsoft~Edge та Apple~Safari (включно з мобільним Safari на iOS). \sked{K45}

\subsection{Процес розробки}

\textbf{REQ-NF-009.} Після кожного запуску тестів агент-розробник зобов'язаний записати log-файл у директорію \texttt{logs/}. Log-файл повинен містити: загальну кількість тестів, кількість passed~/ failed~/ skipped, статус кожного окремого тесту та час виконання тестового набору. \sked{K54}

\subsection{Валідація вхідних даних}

\textbf{REQ-NF-006.} При реєстрації нового користувача система повинна перевіряти:
\begin{itemize}
  \item email~--- формат відповідно до RFC~5322;
  \item пароль~--- мінімальна довжина 8~символів;
  \item нікнейм~--- від 2 до 50~символів.
\end{itemize}
При порушенні будь-якого з обмежень система виводить конкретне повідомлення про помилку, що вказує на причину відмови. \sked{K46}

\section{Критерії прийняття}
\label{sec:kryterii-pryiniattia}

\begin{longtable}{L{2cm} L{2.5cm} L{9.5cm}}
\caption{Критерії прийняття вимог} \label{tab:ac} \\
\toprule
\textbf{AC ID} & \textbf{Вимога} & \textbf{Критерій прийняття} \\
\midrule
\endfirsthead
\toprule
\textbf{AC ID} & \textbf{Вимога} & \textbf{Критерій прийняття} \\
\midrule
\endhead
\midrule
\multicolumn{3}{r}{\emph{Продовження на наступній сторінці}} \\
\endfoot
\bottomrule
\endlastfoot

AC-001 & REQ-F-001 & Після заповнення коректних email, пароля та нікнейму і підтвердження реєстрації користувач отримує доступ до порожнього особистого каталогу. Повторна реєстрація з тим самим email неможлива: система відображає повідомлення про існуючий акаунт. \sked{K7} \\
\addlinespace
AC-002 & REQ-F-002 & Після введення правильних email і пароля користувач перенаправляється до свого каталогу. При неправильних даних~--- повідомлення про помилку, вхід не виконується. \sked{K8} \\
\addlinespace
AC-003 & REQ-F-003 & Після виходу звернення до захищених сторінок перенаправляє на сторінку входу. Повторна авторизація відновлює доступ. \sked{K9} \\
\addlinespace
AC-004 & REQ-F-004 & Каталог відображає всі додані елементи; для кожного видно назву, тип, статус і оцінку. \sked{K10} \\
\addlinespace
AC-005 & REQ-F-005 & При виборі типу «Книги» відображаються лише елементи типу «книга»; при скиданні фільтра~--- знову всі. \sked{K11} \\
\addlinespace
AC-006 & REQ-F-006 & При виборі статусу «завершено» відображаються лише елементи з цим статусом; при скиданні~--- знову всі. \sked{K12} \\
\addlinespace
AC-007 & REQ-F-007 & Після заповнення обов'язкових полів (назва, тип) елемент з'являється в каталозі з відповідними атрибутами та обкладинкою. \sked{K13} \\
\addlinespace
AC-008 & REQ-F-008 & Після редагування та збереження змінені атрибути відображаються в каталозі й на сторінці елемента. \sked{K14} \\
\addlinespace
AC-009 & REQ-F-009 & Після підтвердження видалення елемент зникає з каталогу та з усіх добірок, до яких входив. \sked{K15} \\
\addlinespace
AC-010 & REQ-F-010 & Доданий тег відображається на картці. При введенні в поле тегу~--- підказки раніше використаних тегів. Видалення тегу не впливає на інші атрибути. \sked{K16} \\
\addlinespace
AC-011 & REQ-F-011 & При фільтруванні за тегом відображаються лише позначені ним елементи; результат оновлюється при зміні або скиданні фільтра. \sked{K17} \\
\addlinespace
AC-012 & REQ-F-012 & Для елемента з двома і більше Sessions~--- список усіх записів з датою, оцінкою та нотаткою. \sked{K18} \\
\addlinespace
AC-013 & REQ-F-013 & Sessions відсортовані від найновішої до найстарішої; загальна кількість вказана. \sked{K19} \\
\addlinespace
AC-014 & REQ-F-014 & Після редагування Session~--- оновлені дані. Після видалення~--- сесія зникає, кількість зменшується. \sked{K20} \\
\addlinespace
AC-015 & REQ-F-015 & Елемент одночасно відображається у кожній добірці, до якої його додано. Видалення з однієї не впливає на інші. \sked{K21} \\
\addlinespace
AC-016 & REQ-F-016 & Після введення назви нова порожня добірка з'являється в списку добірок. \sked{K22} \\
\addlinespace
AC-017 & REQ-F-017 & Після збереження нова назва відображається в списку і в заголовку сторінки. \sked{K23} \\
\addlinespace
AC-018 & REQ-F-018 & Після підтвердження видалення добірка зникає; елементи залишаються в каталозі. \sked{K24} \\
\addlinespace
AC-019 & REQ-F-019 & Після видалення елемента з добірки він залишається в каталозі та в інших добірках. \sked{K25} \\
\addlinespace
AC-020 & REQ-F-020 & Сторінка добірки відображає всі елементи з назвою, типом і статусом. \sked{K26} \\
\addlinespace
AC-021 & REQ-F-021 & Якщо дату не заповнено при зміні статусу на «завершено»~--- автоматично підставляється поточна дата. Якщо введено вручну~--- зберігається введене. \sked{K27} \\
\addlinespace
AC-022 & REQ-F-022 & Сторінка статистики відображає кількість завершених елементів за типами та за обраний проміжок. \sked{K28} \\
\addlinespace
AC-023 & REQ-F-023 & При введенні частини назви каталог відображає лише відповідні елементи; при очищенні~--- всі. \sked{K29} \\
\addlinespace
AC-024 & REQ-F-024 & Форма відображає тип-специфічні поля лише після вибору типу; усі поля необов'язкові. \sked{K30} \\
\addlinespace
AC-025 & REQ-F-025 & Без власного зображення~--- дефолтна ілюстрація за типом. Після завантаження~--- власне зображення; його можна видалити, повернувши дефолтне. \sked{K31} \\
\addlinespace
AC-026 & REQ-F-026 & Після вибору параметра сортування каталог перебудовується; при зміні~--- знову. \sked{K32} \\
\addlinespace
AC-027 & REQ-F-027 & Після вибору параметра сортування всередині добірки список перебудовується. \sked{K33} \\
\addlinespace
AC-028 & REQ-F-028 & Спроба завантажити файл невідповідного формату або розміром понад 5~МБ~--- повідомлення про помилку, файл не зберігається. \sked{K48} \\
\addlinespace
AC-029 & REQ-F-029 & Після зміни нікнейму в налаштуваннях акаунту нове ім'я відображається в інтерфейсі. \sked{K49} \\
\addlinespace
AC-030 & REQ-F-022a & Сторінка статистики містить перемикач часового проміжку: «цей місяць», «цей рік», «весь час». \sked{K51} \\
\addlinespace
AC-031 & REQ-NF-006 & При введенні пароля коротше 8~символів~--- конкретне повідомлення. При некоректному email або нікнеймі поза межами 2–50~символів~--- відповідне повідомлення. \sked{K46} \\
\addlinespace
AC-032 & REQ-NF-008 & При порожньому каталозі, порожній добірці або відсутніх результатах пошуку відображається інформативне повідомлення (не порожній список). \sked{K52} \\
\addlinespace
AC-033 & REQ-NF-009 & Після запуску тестів у директорії \texttt{logs/} існує файл з результатами останнього запуску. Файл містить: загальну кількість тестів, кількість passed~/ failed~/ skipped, статус кожного тесту (назва~+ результат), час виконання набору. \sked{K54} \\

\end{longtable}

\section{BDD-сценарії}
\label{sec:bdd}

Нижче наведено BDD-сценарії у форматі Gherkin для ключових функцій системи.

\begin{codelisting}
\caption{BDD-сценарії системи}
\label{lst:bdd}
\begin{skedcode}
Сценарій: Реєстрація нового користувача
  Given відкрита сторінка реєстрації
  And email "test@example.com" ще не зареєстрований у системі
  When користувач вводить email "test@example.com", пароль та нікнейм
    і підтверджує реєстрацію
  Then користувач отримує доступ до порожнього особистого каталогу
  And повторна спроба реєстрації з тим самим email показує повідомлення
    "Акаунт із цим email вже існує"
\end{skedcode}
\end{codelisting}
\sked{K7}

\begin{codelisting}
\caption{Вхід до акаунту}
\label{lst:bdd-login}
\begin{skedcode}
Сценарій: Вхід до акаунту з коректними даними
  Given користувач зареєстрований у системі
  When він вводить правильний email і пароль та натискає "Увійти"
  Then він перенаправляється до свого особистого каталогу
  And бачить свої раніше додані елементи
\end{skedcode}
\end{codelisting}
\sked{K8}

\begin{codelisting}
\caption{Повторний перегляд з іншою оцінкою}
\label{lst:bdd-session}
\begin{skedcode}
Сценарій: Повторний перегляд фільму з іншою оцінкою
  Given елемент контенту "Фільм X" вже має одну Session з оцінкою 9
  When користувач додає нову Session для того самого елемента
    з оцінкою 5 і нотаткою "на другий раз не сподобалось"
  Then елемент контенту має дві Sessions, кожна зі своєю датою,
    оцінкою та нотаткою
  And обидві оцінки зберігаються окремо, жодна не перезаписується
\end{skedcode}
\end{codelisting}
\sked{K18}

\begin{codelisting}
\caption{Елемент у кількох добірках}
\label{lst:bdd-collection}
\begin{skedcode}
Сценарій: Елемент належить кільком добіркам одночасно
  Given користувач має елемент "Книга Y" у каталозі
  And існують добірки "Улюблене" та "2026 рік"
  When користувач додає "Книга Y" до обох добірок
  Then "Книга Y" відображається у вмісті обох добірок
  And видалення "Книга Y" з добірки "2026 рік" не видаляє її
    з "Улюблене" чи з каталогу
\end{skedcode}
\end{codelisting}
\sked{K21,K25}

\begin{codelisting}
\caption{Автозаповнення дати завершення}
\label{lst:bdd-date}
\begin{skedcode}
Сценарій: Автозаповнення дати завершення
  Given користувач редагує елемент зі статусом "у процесі"
  When користувач змінює статус на "завершено"
    і не вказує дату вручну
  Then система автоматично встановлює дату завершення рівною
    поточній даті
\end{skedcode}
\end{codelisting}
\sked{K27}

\begin{codelisting}
\caption{Фільтрування за тегом}
\label{lst:bdd-tag}
\begin{skedcode}
Сценарій: Фільтрування каталогу за тегом
  Given у каталозі є три елементи: два з тегом "фентезі" і один без
  When користувач обирає фільтр за тегом "фентезі"
  Then відображаються лише два елементи, позначені цим тегом
  And при скиданні фільтра знову відображаються всі три елементи
\end{skedcode}
\end{codelisting}
\sked{K17}

\section{Матриця простежуваності вимог}
\label{sec:matrytsia}

\begin{longtable}{L{2.8cm} L{3.2cm} L{2.5cm} L{5.5cm}}
\caption{Матриця простежуваності} \label{tab:matrytsia} \\
\toprule
\textbf{REQ ID} & \textbf{AC ID} & \textbf{BDD} & \textbf{Примітка} \\
\midrule
\endfirsthead
\toprule
\textbf{REQ ID} & \textbf{AC ID} & \textbf{BDD} & \textbf{Примітка} \\
\midrule
\endhead
\midrule
\multicolumn{4}{r}{\emph{Продовження на наступній сторінці}} \\
\endfoot
\bottomrule
\endlastfoot

REQ-F-001 & AC-001 & \lstref{lst:bdd} & Реєстрація \\
REQ-F-002 & AC-002 & \lstref{lst:bdd-login} & Вхід \\
REQ-F-003 & AC-003 & --- & Вихід \\
REQ-F-004 & AC-004 & --- & Каталог \\
REQ-F-005 & AC-005 & --- & Фільтр за типом \\
REQ-F-006 & AC-006 & --- & Фільтр за статусом \\
REQ-F-007 & AC-007 & --- & Додавання елемента \\
REQ-F-008 & AC-008 & --- & Редагування елемента \\
REQ-F-009 & AC-009 & --- & Видалення елемента \\
REQ-F-010 & AC-010 & --- & Теги з автодоповненням \\
REQ-F-011 & AC-011 & \lstref{lst:bdd-tag} & Фільтр за тегом \\
REQ-F-012 & AC-012 & \lstref{lst:bdd-session} & Sessions \\
REQ-F-013 & AC-013 & --- & Перегляд Sessions \\
REQ-F-014 & AC-014 & --- & Редагування~/ видалення Session \\
REQ-F-015 & AC-015 & \lstref{lst:bdd-collection} & Many-to-many добірки \\
REQ-F-016 & AC-016 & --- & Створення добірки \\
REQ-F-017 & AC-017 & --- & Редагування добірки \\
REQ-F-018 & AC-018 & --- & Видалення добірки \\
REQ-F-019 & AC-019 & \lstref{lst:bdd-collection} & Видалення з добірки \\
REQ-F-020 & AC-020 & --- & Перегляд добірки \\
REQ-F-021 & AC-021 & \lstref{lst:bdd-date} & Автодата \\
REQ-F-022 & AC-022 & --- & Статистика \\
REQ-F-022a & AC-030 & --- & Часовий проміжок статистики \\
REQ-F-023 & AC-023 & --- & Пошук за назвою \\
REQ-F-024 & AC-024 & --- & Тип-специфічні поля \\
REQ-F-025 & AC-025 & --- & Обкладинка \\
REQ-F-026 & AC-026 & --- & Сортування каталогу \\
REQ-F-027 & AC-027 & --- & Сортування добірки \\
REQ-F-028 & AC-028 & --- & Обмеження обкладинки \\
REQ-F-029 & AC-029 & --- & Редагування профілю \\
REQ-NF-001 & --- & --- & Продуктивність \\
REQ-NF-002 & --- & --- & Безпека~/ хешування \\
REQ-NF-003 & --- & --- & Надійність \\
REQ-NF-004 & --- & --- & Адаптивний інтерфейс \\
REQ-NF-005 & --- & --- & Сумісність браузерів \\
REQ-NF-006 & AC-031 & --- & Валідація вхідних даних \\
REQ-NF-007 & --- & --- & Безпека сесії (CSRF, HTTP-only) \\
REQ-NF-008 & AC-032 & --- & Порожні стани \\
REQ-NF-009 & AC-033 & --- & Log-файли тестування \\

\end{longtable}


\skedappendix

\section{Функції поза межами MVP}
\label{sec:poza-mvp}

Наступні функції свідомо виключені з версії~v1 з метою стримання обсягу першої реалізації. Вони можуть бути розглянуті у наступних версіях системи. \sked{K43,K50}

\begin{itemize}
  \item Офлайн-режим~--- повноцінна робота без інтернет-з'єднання.
  \item Push- та email-сповіщення.
  \item Багатомовність інтерфейсу.
  \item Функція «друзі»~--- взаємне підписання, спільний доступ до добірок.
  \item Блок рекомендацій або «Популярне серед користувачів».
  \item Інтеграція із зовнішніми базами медіа-контенту (TMDB, Google Books тощо).
  \item Адміністративна панель.
  \item Email-підтвердження при реєстрації (відправка коду на пошту).
  \item Зміна email та зміна пароля в налаштуваннях акаунту.
  \item Видалення акаунту користувачем.
\end{itemize}


\skedbibliography

\end{document}
